\documentclass[11pt]{article}
\usepackage[margin=1in]{geometry}
\usepackage{amsmath, amssymb, amsthm}

\title{Homework 8}
\author{Jonathan Nguyen}
\date{}

\begin{document}
\maketitle

\section*{19.4}
Find all solutions of \(x^2+2x+4=0\) in \(\mathbb{Z}_{6}\).

\noindent\textbf{Solution.}
We test each element of \(\mathbb{Z}_{6}\):
\[
\begin{aligned}
x=0: &\quad 0+0+4=4 \equiv 4 \pmod{6},\\
x=1: &\quad 1+2+4=7 \equiv 1 \pmod{6},\\
x=2: &\quad 4+4+4=12 \equiv 0 \pmod{6},\\
x=3: &\quad 9+6+4=19 \equiv 1 \pmod{6},\\
x=4: &\quad 16+8+4=28 \equiv 4 \pmod{6},\\
x=5: &\quad 25+10+4=39 \equiv 3 \pmod{6}.
\end{aligned}
\]
Therefore, the only solution in \(\mathbb{Z}_{6}\) is \(\boxed{x=2}.\)

\section*{19.10}
Find the characteristic of the ring
\(\mathbb{Z}_{6}\times\mathbb{Z}_{15}\).

\noindent\textbf{Solution.}
The ring \(\mathbb{Z}_{6}\times\mathbb{Z}_{15}\) has unity \((1,1)\).
For a ring with unity, the characteristic is the least positive integer
\(n\) such that \(n\cdot(\text{unity})=0\), so we need the least
positive \(n\) with \(n(1,1)=(0,0)\). Since
\[
n(1,1)=(n\cdot1,\;n\cdot1),
\]
this holds exactly when \(n\cdot1=0\) in \(\mathbb{Z}_6\) and
\(n\cdot1=0\) in \(\mathbb{Z}_{15}\), that is, when \(6\mid n\) and
\(15\mid n\). The least such \(n\) is
\[
\boxed{
\operatorname{char}(\mathbb{Z}_{6}\times\mathbb{Z}_{15})
=\operatorname{lcm}(6,15)=30.
}
\]

\section*{19.23}
An element \(a\) of a ring \(R\) is idempotent if \(a^2=a\).
Show that a division ring contains exactly two idempotent elements.

\noindent\textbf{Solution.}
Let \(R\) be a division ring. Both \(0_R\) and \(1_R\) are idempotent,
since \(0_R^2=0_R\) and \(1_R^2=1_R\).

Now suppose \(a\in R\) is idempotent and \(a\neq0_R\). Since \(R\) is a
division ring, \(a\) has a multiplicative inverse \(a^{-1}\).
Multiplying both sides of \(a^2=a\) on the left by \(a^{-1}\) gives
\[
\begin{aligned}
a^{-1}(aa)&=a^{-1}a,\\
(a^{-1}a)a&=1_R,\\
1_R\,a&=1_R,\\
a&=1_R.
\end{aligned}
\]
Thus every idempotent element of \(R\) is either \(0_R\) or \(1_R\).
Since \(0_R\neq1_R\) in a division ring, these are two distinct elements,
so \(R\) contains exactly two idempotent elements.

\section*{19.26}
Let \(R\) be a ring that contains at least two elements. Suppose that
for each nonzero \(a\in R\), there exists a unique \(b\in R\) such that
\(aba=a\).

\subsection*{(a) Show that \(R\) has no divisors of \(0\).}
\noindent\textbf{Solution.}
Let \(a\in R\) be nonzero, and suppose \(ac=0\) for some \(c\in R\).
Let \(b\) be the unique element of \(R\) such that \(aba=a\).
Since \(aca=(ac)a=0\cdot a=0\),
\[
a(b+c)a=aba+aca=a+0=a.
\]
So \(b+c\) also satisfies the defining equation for \(a\).
By uniqueness, \(b+c=b\), and hence \(c=0\).

We have shown that if \(a\neq0\) and \(ac=0\), then \(c=0\). So there
are no nonzero \(a,c\in R\) with \(ac=0\), which means \(R\) has no
divisors of \(0\). Equivalently, if \(xy=0\) in \(R\), then \(x=0\) or
\(y=0\).

\subsection*{(b) Show that \(bab=b\).}
\noindent\textbf{Solution.}
Let \(a\neq0\), and let \(b\) be the unique element with \(aba=a\).
Multiplying \(aba=a\) on the right by \(b\) gives \(abab=ab\).
Hence
\[
a(bab-b)=abab-ab=ab-ab=0.
\]
Since \(a\neq0\) and \(R\) has no divisors of \(0\) by part (a),
\(bab-b=0\). Therefore,
\[
\boxed{bab=b.}
\]

\subsection*{(c) Show that \(R\) has unity.}
\noindent\textbf{Solution.}
Since \(R\) contains at least two elements, there is a nonzero
\(a\in R\). Let \(b\) be the unique element with \(aba=a\), and set
\(e=ab\). Then
\[
ea=(ab)a=aba=a\neq0,
\]
so \(e\neq0\) (if \(e=0\), then \(ea=0\)). Also,
\[
e^2=(ab)(ab)=(aba)b=ab=e.
\]
Now let \(x\in R\) be arbitrary. Using \(e^2=e\),
\[
e(ex-x)=e^2x-ex=ex-ex=0.
\]
Since \(e\neq0\) and \(R\) has no divisors of \(0\) by part (a),
\(ex-x=0\), so \(ex=x\). Similarly,
\[
(xe-x)e=xe^2-xe=xe-xe=0,
\]
and since \(e\neq0\), \(xe-x=0\), so \(xe=x\).
Thus \(ex=xe=x\) for every \(x\in R\), so \(e\) is a unity of \(R\).

\subsection*{(d) Show that \(R\) is a division ring.}
\noindent\textbf{Solution.}
By part (c), \(R\) has unity, which we denote by \(1_R\).
Since \(R\) has at least two elements, \(1_R\neq0\): otherwise every
\(x\in R\) would satisfy \(x=x\cdot1_R=x\cdot0=0\), so \(R=\{0\}\).

Now let \(a\in R\) be any nonzero element, and let \(b\) be the unique
element with \(aba=a\). Then
\[
a(ba-1_R)=aba-a=0
\quad\text{and}\quad
(ab-1_R)a=aba-a=0.
\]
Since \(a\neq0\) and \(R\) has no divisors of \(0\) by part (a),
\(ba-1_R=0\) and \(ab-1_R=0\), so
\[
ab=ba=1_R.
\]
Thus \(b=a^{-1}\), and every nonzero element of \(R\) has a
multiplicative inverse. Since \(R\) is a ring with unity \(1_R\neq0\)
in which every nonzero element has a multiplicative inverse, \(R\) is a
division ring.

\section*{20.2}
Find a generator of the multiplicative group of nonzero elements
of \(\mathbb{Z}_{11}\).

\noindent\textbf{Solution.}
Since \(11\) is prime, \(\mathbb{Z}_{11}\) is a field, so its nonzero
elements form a group \(\mathbb{Z}_{11}\setminus\{0\}\) under
multiplication, and this group has order \(10\).
By Lagrange's theorem, the order of \(2\) in this group divides \(10\),
so it is \(1,2,5,\) or \(10\). We compute
\[
2^2=4\not\equiv1\pmod{11}
\quad\text{and}\quad
2^5=32=2(11)+10\equiv10\not\equiv1\pmod{11}.
\]
If \(2\) had order \(1\) or \(2\), then \(2^2\equiv1\pmod{11}\); if it had
order \(5\), then \(2^5\equiv1\pmod{11}\). Neither holds, so \(2\) has
order \(10\). Since the group has order \(10\), the cyclic subgroup
generated by \(2\) is the whole group. Therefore, \(2\) generates
\(\mathbb{Z}_{11}\setminus\{0\}\):
\[
\boxed{\langle 2\rangle=\mathbb{Z}_{11}\setminus\{0\}.}
\]

\section*{20.4}
Use Fermat's theorem to find the remainder of \(3^{47}\)
when it is divided by \(23\).

\noindent\textbf{Solution.}
Since \(23\) is prime and \(23\nmid3\), Fermat's little theorem gives
\[
3^{22}\equiv1\pmod{23}.
\]
Since \(47=2(22)+3\),
\[
3^{47}=(3^{22})^2\cdot 3^3
\equiv1^2\cdot27
=27
\equiv4\pmod{23},
\]
because \(27=23+4\). Since \(0\le4<23\), the remainder is
\(\boxed{4}\).

\section*{20.12}
Describe all the solutions of \(22x\equiv5\pmod{15}\).

\noindent\textbf{Solution.}
Since \(22=15+7\), we have \(22\equiv7\pmod{15}\), so the congruence
becomes
\[
7x\equiv5\pmod{15}.
\]
Since \(\gcd(7,15)=1\), \(7\) has a multiplicative inverse modulo
\(15\). Indeed, \(7(13)=91=6(15)+1\equiv1\pmod{15}\).
Multiplying both sides by \(13\) gives
\[
\begin{aligned}
13(7x)&\equiv13(5)\pmod{15},\\
91x&\equiv65\pmod{15},\\
x&\equiv65\pmod{15},\\
x&\equiv5\pmod{15},
\end{aligned}
\]
using \(91\equiv1\) and \(65=4(15)+5\). Multiplying by \(13\) can be
undone by multiplying by \(7\), so each step is reversible and these are
exactly the solutions. As a check, \(22(5)=110=7(15)+5\equiv5\pmod{15}\).
Thus, all integer solutions are
\[
\boxed{x=5+15k,\qquad k\in\mathbb{Z}.}
\]

\section*{20.27}
Show that $1$ and $p - 1$ are the only elements of the field
$\mathbb{Z}_p$ that are their own multiplicative inverse.

\noindent\textbf{Solution.}
An element $x \in \mathbb{Z}_p$ is its own multiplicative inverse exactly
when $x \cdot x = 1$, that is, when $x^2 = 1$. Thus
\begin{align*}
    x^2 - 1 &= (x + 1)(x - 1) = 0.
\end{align*}
Since $\mathbb{Z}_p$ is a field, it is an integral domain and has no
divisors of zero, so $x + 1 = 0$ or $x - 1 = 0$.
In $\mathbb{Z}_p$, this means $x = -1 = p - 1$ or $x = 1$.

Conversely, both of these elements are their own inverses:
$1 \cdot 1 = 1$, and
\begin{align*}
    (p - 1)^2 &= p^2 - 2p + 1 \equiv 1 \pmod{p}.
\end{align*}
Therefore, $1$ and $p - 1$ are the only elements of $\mathbb{Z}_p$ that
are their own multiplicative inverse. (When $p = 2$, these are the same
element, since $p - 1 = 1$.)

\end{document}